\section{The Numbers Behind the Reassurance}

\subsection{The provenance rule this section obeys}

\draftinstr{Roughly 900 characters. Every quantitative claim in this section
carries one of four class letters, and the letter appears on the same line as
the number. M means measured in a cited human series; C means a contemporaneous
nationwide comparator; S means an author simulation and not human data; P means
a protocol-specified limit and not a result. State that two of the nine headline
values are class S and name them. Cite \texttt{Jauniaux2025} for why this
discipline matters: expectation and experience diverge systematically in robotic
surgery, and expectation is usually set by marketing.}

\subsection{The reassurance dashboard}

\draftinstr{Roughly 1,200 characters introducing Figure 27 panel by panel, and
do not skip the unfavourable panels. Panel A: fistula and mortality falling with
institutional volume, cite \texttt{DutchCohort2025}. Panel B: 8.2 percent of the
1000 Phase 0 simulated procedures halted by a safety mechanism, cite
\texttt{pdac060s2030}. Panel C: every observed worst case inside its cap. Panel
D: the survival context with the honest denominator, cite \texttt{Siegel2025}
and \texttt{Conroy2018}. Panel E: how many concerns get which class of answer.}

\begin{pafig}[mermaid-type: xychart + pie + CI strip, full page]
\node[mmgoal] (a) at (0,0) {Figure 27 slot\\five-panel dashboard};
\node[mmin] (b) at (0,-2.0) {\ttfamily\tiny Source\\ mermaid/\\ fig-\pfb 27-\pfb reassurance-\pfb dashboard.md};
\draw[mmedge] (a) -- (b);
\end{pafig}
\vspace{-0.7cm}
\figcaption{Figure 27. Five panels of quantitative reassurance on one page, learning-curve
outcomes, Phase 0 simulation endings, safety-limit headroom, the survival context with\\
its honest denominator, and the class of answer each documented concern receives.}

\draftinstr{Replace the Figure 27 slot with the full-page mermaid-type dashboard
from \texttt{mermaid/\pfb fig-\pfb 27-\pfb reassurance-\pfb dashboard.md}. Five stacked panels
inside one frame separated by 0.4pt \texttt{pagraym} hairlines; Panel A with
\texttt{\textbackslash vbarcol} plus a 0.8pt polyline; Panel B with
\texttt{\textbackslash donutseg} and leader lines to labels outside the ring;
Panel C with \texttt{\textbackslash vbarpair}; Panels D and E as
\texttt{tabularx} tables at the panel measure.}

\subsection{Robotic pancreatoduodenectomy, by learning-curve phase}

\draftinstr{Populate from \texttt{DutchCohort2025}. Add the class letter to
every numeric column heading. State the protocol's response in the final column:
enrollment is restricted to a site already past phase 3, so the learning curve
is not paid for by the participant.}

\vspace{0.30em}

\noindent\begin{tabularx}{\textwidth}{L{3.6cm} C{2.5cm} C{2.4cm} Y}
\toprule
\textbf{Learning-curve phase} & \textbf{ISGPS B/C fistula (C)} & \textbf{90-day mortality (C)} & \textbf{What this protocol does about it} \\
\midrule
Phase 1, cases 1 to 80 & 24.0 percent & 8.0 percent & Not eligible to enrol; the site must be past this phase \\
Phase 2, cases 81 to 180 & 18.5 percent & 6.0 percent & Not eligible to enrol \\
Phase 3, cases 181 to 400 & 14.0 percent & 4.5 percent & Minimum eligible experience for a participating site \\
Phase 4, cases 401 and above & 11.0 percent & 3.5 percent & Preferred; operator case counts are published to participants \\
\bottomrule
\end{tabularx}

\vspace{0.30em}

\subsection{Safety-limit headroom}

\draftinstr{Populate from \texttt{onpremwhippl} for the caps and
\texttt{pdac060s2030} for the observed values. Mark the cap column P and the
observed column S, and repeat in the caption that a simulated worst case is not
a human worst case.}

\vspace{0.30em}

\noindent\begin{tabularx}{\textwidth}{L{4.4cm} C{2.4cm} C{2.4cm} Y}
\toprule
\textbf{Limit} & \textbf{Protocol cap (P)} & \textbf{Worst observed (S)} & \textbf{Headroom} \\
\midrule
Tip force, any single arm & 3.0 N & 2.1 N & 30 percent \\
Tip force, all arms summed & 18.0 N & 12.4 N & 31 percent \\
Cross-arm emergency stop & 3.0 ms & 1.7 ms & 43 percent \\
System-wide emergency stop & 500 ms & 310 ms & 38 percent \\
Positional error & 2.0 mm & 1.3 mm & 35 percent \\
Force reproducibility & 0.5 N & 0.31 N & 38 percent \\
\bottomrule
\end{tabularx}

\vspace{0.30em}

\subsection{The survival context, with the honest denominator}

\draftinstr{Roughly 1,000 characters and the most important honest passage in
the paper. Pancreatic cancer five-year overall survival across all stages is
13 percent \cite{Siegel2025}. Resected PDAC after adjuvant FOLFIRINOX is
materially better \cite{Conroy2018}. An R0 margin is materially better than R1.
A Phase 1 study of at most eighteen participants cannot move a survival curve
and this paper does not claim it will. What it can establish is whether the
combination is safe enough to be tested in a study that could.}

\subsection{Where the numbers come from}

\draftinstr{Roughly 800 characters introducing Figure 28. Nine quoted values,
six sources, four evidence classes, and one rule: a value with no incoming edge
does not appear in this paper.}

\begin{pafig}[graphviz-type: provenance DAG, record nodes]
\node[gvkey] (a) at (0,0) {Figure 28 slot\\evidence provenance};
\node[gvbox] (b) at (0,-2.0) {\ttfamily\tiny Source\\ graphviz/\\ fig-\pfb 28-\pfb evidence-\pfb provenance.dot};
\draw[gvedge] (a) -- (b);
\end{pafig}
\vspace{-0.7cm}
\figcaption{Figure 28. Every quantitative claim in this paper traced through its source to
its class of evidence, with the two class-S values, which are simulated rather than\\
measured, marked by dashed edges so a reader can separate them at a glance.}

\draftinstr{Replace the Figure 28 slot with the graphviz-type provenance DAG
from \texttt{graphviz/\pfb fig-\pfb 28-\pfb evidence-\pfb provenance.dot}. Three ranks at
$x = 0, 6.4, 12.4$; nine \texttt{gvrec} record nodes on rank 1 at 1.9 cm pitch;
six \texttt{gvbox} sources on rank 2; four \texttt{gvnode} class ellipses on
rank 3; class-S edges dashed \texttt{pagrayd}; where three or more values share
a source, fan the outer edges at looseness 0.8.}

\subsection{What would change the conclusion}

\draftinstr{Roughly 700 characters. Name the findings that would make this paper
wrong: a device-related serious adverse event rate above the prespecified halt
threshold, an R0 rate below the comparator range, a fistula rate above phase-1
learning-curve levels at an experienced site, or any emergency-stop latency
measured above budget in a human procedure. State that each has a prespecified
stopping rule attached, and forward to \S5.6.}
